\section{Introduction}

Biological diversity underpins the ecosystem services on which human societies depend, yet global assessments report that those services are declining at historically unprecedented rates \citep{ipbes2019global,newbold2015land}. In response, the Kunming-Montreal Global Biodiversity Framework adopted a portfolio of commitments intended to reverse nature loss by 2030, including the designation of thirty percent of the Earth's land and ocean as protected areas (30by30). Reaching this target requires knowledge of where species occur at spatiotemporal resolutions fine enough to support conservation planning, adaptive management, and the emerging practice of nature-related financial disclosure. Such knowledge has traditionally been produced through on-site surveys conducted by specialists, but expert effort alone cannot scale to the resolution and geographic coverage now demanded.

Over the past two decades, community-sourced observations collected through web and mobile platforms such as eBird and iNaturalist have begun to fill this gap \citep{sullivan2009ebird,bonney2009citizen,chandler2017contribution}. When aggregated at large scales, these records have supported detection of phenological shifts, tracking of invasive species, mapping of trait variation, and species distribution modelling \citep{chandler2017contribution,pocock2015brc}. At the same time, such records are typically collected opportunistically and suffer from well-documented spatial, taxonomic, and observer biases \citep{isaac2014statistics,feldman2021trends,kendal2020citybias}, which complicates statistical inference and constrains their utility for protected-area prioritisation.

Two complementary strategies can improve the information yield from community-sourced biodiversity data. The first broadens participation by embedding biological recording into applications that are pleasurable to use, often through gamification elements that reward submissions and cross-user interaction \citep{bonney2009citizen}. The second tightens the statistical treatment of presence-only records by explicitly modelling sampling effort, most prominently through maximum-entropy species distribution models (MaxEnt) and their bias-correction extensions \citep{phillips2006maxent,phillips2008maxentextensions,phillips2009samplebias,elith2011statistical}. Prior work has shown that, when sampling effort is correctly handled, pooling community observations with structured surveys produces more accurate distribution estimates than either source in isolation \citep{johnston2018estimates,johnston2021analytical,robinson2020integrating,steen2019stringent,milanesi2020observer}. Almost all of this evidence, however, comes from a narrow set of well-recorded taxa, with birds dominating the published record \citep{feldman2021trends}.

In this work we ask whether a new generation of smartphone-first, AI-assisted, gamified biodiversity applications can deliver community-sourced records of sufficient quality and environmental coverage to meaningfully improve species distribution models across a broad taxonomic spread, including taxa such as seed plants that have proven difficult in prior evaluations. We focus on Biome, a Japan-deployed mobile application launched in 2019, and pair an audit of its observational records with a controlled comparison of MaxEnt models trained on traditional survey data alone versus a blend that includes Biome records. Our contributions are:

\begin{itemize}
  \item A large-scale data-quality audit of 1,420 Biome records spanning eight taxonomic groups, reporting species-, genus-, and family-level identification accuracy stratified by wild-individual status and species rarity.
  \item A benchmark comparison of MaxEnt species distribution models trained on Traditional versus Biome+Traditional data for 132 terrestrial species, using a 25\%--75\% test mixture designed to avoid favouring either training regime and a spatially blocked evaluation protocol.
  \item Evidence that incorporating Biome records reduces the volume of training data required to reach high SDM accuracy by roughly an order of magnitude, with the largest gains for endangered species.
  \item A discussion of how smartphone-driven community science can concretely support the Kunming-Montreal Global Biodiversity Framework and its implementation pathways, including protected-area designation and corporate nature-related financial disclosure.
\end{itemize}

\section{Related Work}

\paragraph{Community science platforms and data quality.}
Structured and semi-structured platforms such as eBird \citep{sullivan2009ebird} have demonstrated that broad public participation can generate datasets large enough to support macroecological analyses, while the Biological Records Centre and similar long-standing schemes have shown that volunteer-collected data can be curated to research-grade standards over decades \citep{pocock2015brc,bonney2009citizen,chandler2017contribution}. Independent audits of identification accuracy are increasingly available: evaluations of plant identification apps report top-one accuracy well below that of structured expert review, with the best free tools reaching the mid-eighties and popular mobile apps in the sixty to seventy percent range \citep{hart2023plantapps}, while early volunteer programs for invasive plants achieved roughly seventy percent correctness \citep{crall2011assessing}. These baselines motivate our explicit per-taxon accuracy evaluation of Biome records.

\paragraph{Presence-only species distribution modelling.}
MaxEnt is widely used for presence-only SDMs because it balances predictive accuracy with computational efficiency and supports principled treatment of sampling bias via target-group background selection and informative pseudo-absence schemes \citep{phillips2006maxent,phillips2008maxentextensions,phillips2009samplebias,elith2011statistical}. Recent benchmark work confirms that MaxEnt remains competitive against more flexible machine-learning alternatives when predictor sets and regularisation are chosen carefully \citep{valavi2022predictive}. The ENMeval 2.0 package \citep{kass2021enmeval} provides automated tuning and evaluation infrastructure for MaxEnt and related ecological niche models, and underlies our pipeline. Model evaluation for presence-only data is typically performed via the Boyce index \citep{hirzel2006boyce}, which measures the correlation between predicted habitat suitability and the observed frequency of presences, and is suited to settings without verified absences.

\paragraph{Integrating community and structured data for SDMs.}
Several studies have shown that blending community-sourced and structured-survey data can improve SDM accuracy, most clearly for birds \citep{johnston2018estimates,johnston2021analytical,robinson2020integrating,steen2019stringent}. Proposed mechanisms include an increase in effective sample size, complementary environmental coverage, and the use of observer-oriented pseudo-absence schemes that leverage sampling effort proxies from the community dataset \citep{milanesi2020observer}. A recent quantitative review highlights two persistent gaps: a strong taxonomic bias toward birds and mammals, and inconsistent attention to spatial and observer biases in published analyses \citep{feldman2021trends}. Our analysis addresses both gaps by evaluating 132 species spanning plants and animals under a unified sampling-effort-aware pseudo-absence scheme, and by comparing Traditional-only versus Biome-blended training across multiple data volumes.

\paragraph{Citizen-science engagement and the social side of data collection.}
The rate at which community-sourced records accumulate is ultimately bounded by participation. Reviews of citizen-science engagement argue that enjoyment, social feedback, and a sense of contribution are primary motivators for sustained participation \citep{bonney2009citizen,chandler2017contribution}, and growing concern about the erosion of everyday contact with nature underscores the indirect value of platforms that encourage outdoor observation \citep{soga2023biophobia,kendal2020citybias}. Biome's design, which combines AI-assisted identification with gamified points and levels, situates it within this broader movement to make biological recording both statistically useful and socially rewarding.
